\section{Conclusiones y recomendaciones al cliente}
El análisis complejo aporta una construcción: pasa de la métrica esférica a una función analítica, identifica $\ln|f'|$ como potencial armónico y determina su frontera mediante Dirichlet. La comparación demuestra cuánto se gana frente al catálogo: la oscilación de escala baja siete veces respecto a Lambert y 76 respecto a Mercator. Proyectar permite conservar algoritmos planos, aunque no reemplaza toda operación geodésica.

El beneficio termina donde cambian las condiciones del modelo: áreas, recorridos fuera del dominio, exigencia subporcentual frente al elipsoide o distancias comparables con los saltos numéricos. En esos casos se requieren proyecciones equivalentes, una región ampliada o cálculo geodésico directo.

Se recomienda EPSG:6372 como sistema de trabajo interoperable. Con un presupuesto ilustrativo de error de 50 m, su error mediano muestral equivale a esa magnitud cerca de 17 km, y su máximo observado cerca de 7 km. La envolvente desplaza el primer umbral a 33 km, pero no mejora sustancialmente el segundo. Estos umbrales son orientativos, no fronteras universales. Mantener la óptima se justifica si esa reducción típica compensa su costo y el cliente valida sus propios datos.
